\chapter{Stato dell'arte}
\label{chap:stato-arte}

Definiti il problema e i criteri di valutazione, occorre esaminare i metodi con cui apprendere la normalità visiva e riconoscerne le deviazioni. Nell'ispezione industriale, l'obiettivo è riconoscere difetti anche quando non sono disponibili esempi rappresentativi di tutte le anomalie possibili \cite{Yang_et_al._2021,Cui_et_al._2023}.

Le applicazioni comprendono il controllo di qualità, l'analisi di immagini mediche e il rilevamento di eventi nei video. Questi ambiti condividono il problema della scarsità delle anomalie, ma differiscono per dati, vincoli e modalità di valutazione. La rassegna seguente privilegia l'ispezione di immagini industriali, che costituisce il dominio della sperimentazione.

\section{Criteri di classificazione dei metodi}

La letteratura può essere organizzata secondo tre assi distinti:
\begin{itemize}
    \item \textbf{dominio dei dati}: immagini, video o rappresentazioni tridimensionali;
    \item \textbf{informazioni disponibili durante l'apprendimento}: campioni normali, esempi di anomalie e annotazioni di diversa granularità;
    \item \textbf{meccanismo di rilevamento}: ricostruzione, distanze nello spazio delle feature, stima di densità o discriminazione appresa.
\end{itemize}

Questi criteri non sono intercambiabili. Un'architettura non determina da sola il paradigma di supervisione e uno stesso metodo può combinare più meccanismi. La tassonomia serve quindi a chiarire il ruolo delle componenti, evitando di considerare le famiglie come categorie necessariamente disgiunte \cite{Cui_et_al._2023}.

\subsection{Granularità dell'output e supervisione}

Il rilevamento a livello di immagine produce un punteggio per l'intera immagine; la localizzazione a livello di pixel produce una mappa spaziale. Una segmentazione binaria può essere ottenuta applicando una soglia alla mappa. L'esistenza di un output spaziale non implica che il modello sia stato addestrato usando maschere dei difetti.

Come definito nel Capitolo~\ref{chap:introduzione}, il progetto adotta uno scenario \textit{one-class}: i campioni industriali utilizzati per l'adattamento sono tutti normali. Le maschere del test sono impiegate nella valutazione. I modelli possono comunque utilizzare rappresentazioni pre-addestrate su altri dataset. Altri metodi one-class introducono anomalie sintetiche per costruire compiti auto-supervisionati; questa procedura non è utilizzata dalle quattro pipeline sperimentate.

La Tabella~\ref{tab:taxonomy} distingue le principali informazioni che possono guidare l'apprendimento.
\begin{table}[htbp]
    \centering
    \small
    \caption{Informazioni di addestramento e ruolo delle annotazioni nell'ispezione visiva.}
    \label{tab:taxonomy}
    \begin{tabularx}{\textwidth}{@{}P{0.24\textwidth}YY@{}}
        \hline
        \textbf{Impostazione} & \textbf{Dati disponibili} & \textbf{Informazione utilizzata} \\
        \hline
        Supervisionata & Campioni normali e difettosi & Etichette di classe o maschere dei difetti \\
        Debolmente supervisionata & Annotazioni meno dettagliate dell'output richiesto & Ad esempio, etichette di immagine per localizzare un difetto \\
        One-class & Campioni industriali noti come normali & Modellazione della normalità senza esempi reali dei difetti \\
        Non supervisionata in senso generale & Campioni non etichettati, eventualmente contaminati & Regolarità statistiche dei dati e ipotesi sulle anomalie \\
        \hline
    \end{tabularx}
\end{table}

\section{Metodi basati su descrittori manuali}

Prima della diffusione delle reti profonde, molte pipeline combinavano descrittori progettati manualmente (\textit{handcrafted features}) con classificatori o stimatori statistici \cite{Yang_et_al._2021}. L'efficacia dipendeva dalla capacità dei descrittori di rappresentare le variazioni rilevanti per il dominio.

Tra gli strumenti utilizzati si distinguono:
\begin{itemize}
    \item \textbf{descrittori locali}: SIFT (\textit{Scale-Invariant Feature Transform}) per punti di interesse invarianti alla scala, HOG (\textit{Histograms of Oriented Gradients}) per orientazioni dei gradienti e LBP (\textit{Local Binary Patterns}) per regolarità locali della texture \cite{Yang_et_al._2021,Cui_et_al._2023};
    \item \textbf{filtri e trasformate}: operatori di gradiente, rilevamento dei bordi, filtri di Gabor e trasformate di Fourier o wavelet per descrivere frequenze, orientazioni e discontinuità \cite{Cui_et_al._2023}.
\end{itemize}

Le rappresentazioni potevano essere utilizzate da tre principali classi di metodi decisionali:
\begin{enumerate}
    \item \textbf{Stima di densità}. Modelli a mistura gaussiana (GMM) e stimatori di densità basati su kernel (KDE) descrivono la distribuzione dei campioni. Una densità bassa rispetto al modello può segnalare un'anomalia.
    \item \textbf{Classificazione one-class}. One-Class SVM (OCSVM) e Support Vector Data Description (SVDD) costruiscono un confine attorno ai campioni normali nello spazio delle feature.
    \item \textbf{Modelli spaziali e temporali}. I campi aleatori di Markov (MRF) e altri modelli contestuali considerano dipendenze tra osservazioni. Nell'analisi video possono essere affiancati da descrittori del movimento \cite{Tran_et_al._2022}.
\end{enumerate}

L'alta dimensionalità e le variazioni di illuminazione o texture possono rendere difficile la scelta delle feature. Inoltre, l'estrazione e la decisione sono spesso ottimizzate separatamente. Questi limiti hanno motivato lo sviluppo di rappresentazioni apprese; i metodi tradizionali rimangono tuttavia utili quando i dati sono limitati o i vincoli computazionali favoriscono soluzioni semplici \cite{Yang_et_al._2021}.

\section{Famiglie di metodi basati sul deep learning}

Le reti profonde permettono di apprendere descrittori multilivello e di combinare informazione locale e rappresentazioni più astratte. Alcuni metodi addestrano una rete sul dominio industriale; altri utilizzano un estrattore pre-addestrato e stimano soltanto un modello delle feature normali.

Le sezioni seguenti presentano cinque famiglie utili a inquadrare i quattro modelli selezionati. La sintesi evidenzia sia i principi condivisi sia le differenze che dovranno essere considerate nel confronto sperimentale \cite{Cui_et_al._2023}.

\subsection{Metodi basati sulla ricostruzione}

I metodi ricostruttivi apprendono a riprodurre immagini normali attraverso una rappresentazione latente. L'ipotesi è che i difetti, non osservati durante l'addestramento, siano ricostruiti meno fedelmente e producano una discrepanza maggiore \cite{Yang_et_al._2021}. Questa proprietà deve essere verificata: una rete sufficientemente flessibile può ricostruire anche anomalie.

Negli autoencoder, l'encoder comprime l'input e il decoder ne genera una ricostruzione. La funzione obiettivo può misurare differenze tra pixel o utilizzare indici strutturali. Una possibile combinazione è
\[
    \mathcal{L}_{\mathrm{rec}}(x,\widehat{x}) =
    (1-\alpha)\lVert x-\widehat{x}\rVert_2
    +\alpha\bigl(1-\mathrm{SSIM}(x,\widehat{x})\bigr),
    \qquad 0\leq\alpha\leq1.
\]
La formula illustra un criterio ricostruttivo e non una perdita universale della famiglia. Il calcolo di una mappa locale richiede inoltre di specificare come le discrepanze siano valutate nello spazio.

Nei modelli generativi come AnoGAN, la ricerca della ricostruzione più vicina può richiedere un'ottimizzazione del vettore latente. Il criterio combina un residuo dell'immagine con una discrepanza tra feature del discriminatore:
\[
    \mathcal{L}_{\mathrm{AnoGAN}}(z)=
    (1-\lambda)\lVert x-G(z)\rVert_1
    +\lambda\lVert D_{\mathrm{feat}}(x)
    -D_{\mathrm{feat}}(G(z))\rVert_1.
\]
Qui $G$ è il generatore e $D_{\mathrm{feat}}$ indica una rappresentazione intermedia del discriminatore \cite{Schlegl_et_al._2017}.

Tra gli sviluppi della famiglia, CAVGA introduce vincoli di attenzione, mentre DFR (\textit{Deep Feature Reconstruction}) ricostruisce rappresentazioni estratte da una rete pre-addestrata anziché direttamente i pixel. La ricostruzione può quindi operare in spazi differenti e sovrapporsi ad altri approcci della tassonomia \cite{Cui_et_al._2023}.

CCW-SSIM, selezionato per il progetto, appartiene a questa famiglia: combina la ricostruzione mediante autoencoder con una valutazione della discrepanza strutturale e cromatica \cite{Berasi_2023}. Il suo interesse per il confronto consiste nell'apprendere direttamente dai campioni industriali senza utilizzare un backbone ImageNet.

\subsection{Metodi basati sulle rappresentazioni}

Questi metodi descrivono i campioni mediante vettori di caratteristiche, indicati come \textit{feature}, spesso estratte da reti pre-addestrate. Il punteggio di anomalia deriva dal confronto con feature normali conservate in memoria oppure con statistiche stimate su di esse \cite{Cui_et_al._2023}.

La famiglia comprende approcci con estrattore congelato e varianti che adattano le rappresentazioni al dominio. Non è quindi corretto attribuire a tutti l'assenza di addestramento neurale.

PaDiM (\textit{Patch Distribution Modeling}) stima una gaussiana multivariata per ciascuna posizione spaziale. Il punteggio locale usa la distanza di Mahalanobis:
\[
    D_M(z)=\sqrt{(z-\mu)^{\mathrm{T}}\Sigma^{-1}(z-\mu)},
\]
dove $z$ è il vettore di feature, $\mu$ la media locale e $\Sigma$ la covarianza stimata sui campioni normali. La distanza considera variabilità e correlazioni tra componenti, a differenza della distanza euclidea \cite{Defard_et_al._2021}.

Altri esempi sono:
\begin{itemize}
    \item \textbf{SPADE}: confronta feature multiscala attraverso ricerche dei vicini più prossimi; la consultazione dei campioni di riferimento contribuisce al costo di inferenza \cite{Yang_et_al._2021};
    \item \textbf{PatchCore}: conserva feature normali in una \textit{memory bank}, riducendola mediante selezione di un sottoinsieme rappresentativo (\textit{coreset}) \cite{Roth_et_al._2022};
    \item \textbf{CFA} (\textit{Coupled-hypersphere-based Feature Adaptation}): adatta le feature al dominio attraverso vincoli geometrici \cite{Cui_et_al._2023}.
\end{itemize}

Il vantaggio dell'estrattore pre-addestrato è la disponibilità di rappresentazioni già strutturate. Il limite è la possibile distanza tra il dominio del pre-addestramento e le caratteristiche dei materiali industriali. Memoria, costo di adattamento e inferenza dipendono poi dalla specifica modalità di confronto.

\subsection{Metodi basati sulla distillazione Teacher--Student}

La distillazione utilizza una rete \textit{Teacher} per fornire rappresentazioni di riferimento e una rete \textit{Student} addestrata a riprodurle sui campioni normali. In inferenza, la discrepanza tra le rappresentazioni può segnalare una deviazione dalla normalità \cite{Cui_et_al._2023}.

Indicando con $f_T^k(x)$ e $f_S^k(x)$ le feature corrispondenti al livello $k$, una forma generale dell'obiettivo è
\[
    \mathcal{L}_{\mathrm{dist}}=
    \sum_k d\bigl(f_T^k(x),f_S^k(x)\bigr),
\]
dove la distanza $d$ e la modalità di aggregazione dipendono dal metodo. Il punteggio spaziale in inferenza può utilizzare un'aggregazione diversa da quella impiegata nell'addestramento.

RD4AD adotta la \textit{reverse distillation}: lo Student è un decoder che riceve una rappresentazione compressa delle feature del Teacher e tenta di ricostruirle. Un \textit{One-Class Bottleneck Embedding} (OCBE) aggrega le informazioni multiscala. La struttura asimmetrica mira a limitare il trasferimento indiscriminato delle informazioni anomale allo Student \cite{Deng_et_al._2022}.

Il metodo ricostruisce dunque feature anziché immagini. Questa relazione con gli approcci ricostruttivi chiarisce perché la tassonomia sia uno strumento descrittivo e non una separazione rigida. Nel progetto, RD4AD permette di confrontare la ricostruzione distillata delle feature con la modellazione statistica di PaDiM.

\subsection{Metodi basati sui flussi di normalizzazione}

I \textit{normalizing flows} apprendono trasformazioni invertibili tra una distribuzione complessa e una distribuzione di riferimento trattabile. Nell'ispezione visiva possono operare sulle feature estratte da un backbone \cite{Cui_et_al._2023}.

Se $z=f_\theta(x)$ è una trasformazione invertibile, il cambio di variabile fornisce
\[
    p_X(x)=p_Z(f_\theta(x))
    \left|\det\frac{\partial f_\theta(x)}{\partial x}\right|,
\]
e quindi
\[
    \log p_X(x)=\log p_Z(f_\theta(x))
    +\log\left|\det\frac{\partial f_\theta(x)}{\partial x}\right|.
\]
La distribuzione $p_Z$ è spesso una normale standard. L'addestramento massimizza la log-verosimiglianza dei campioni normali; il punteggio di anomalia in inferenza deve essere definito in modo coerente con la specifica implementazione. Una densità bassa non equivale direttamente a una probabilità calibrata di difetto.

Tra gli esempi della famiglia:
\begin{itemize}
    \item \textbf{DifferNet} applica flussi alle feature per il rilevamento a livello di immagine;
    \item \textbf{CFLOW-AD} impiega flussi condizionali per la localizzazione;
    \item \textbf{FastFlow} utilizza trasformazioni bidimensionali sulle mappe di feature, mantenendo la struttura spaziale utile alla localizzazione \cite{Yu_et_al._2021}.
\end{itemize}

FastFlow è il rappresentante selezionato nel progetto. La formulazione generale permette estrattori differenti, ma il confronto documentato utilizza un backbone convoluzionale congelato. Il Capitolo~\ref{chap:implementazione} distingue la funzione obiettivo del flusso dalla costruzione effettiva della mappa di anomalia.

\subsection{Metodi basati sulla sintesi di anomalie}

Un'ulteriore famiglia genera alterazioni artificiali dei campioni normali per costruire un segnale di apprendimento. L'assenza di esempi reali dei difetti non impedisce così l'impiego di etichette sintetiche per compiti di classificazione o segmentazione \cite{Cui_et_al._2023}.

Per la segmentazione, una possibile componente dell'obiettivo è l'entropia incrociata binaria:
\[
    \ell(y,p)=-\bigl[y\log p+(1-y)\log(1-p)\bigr],
\]
dove $y$ indica se il pixel appartenga alla regione alterata e $p$ è la previsione della rete. La funzione è calcolata sui pixel e aggregata secondo il metodo; possono essere presenti altre componenti di perdita.

Esempi di questa famiglia sono:
\begin{itemize}
    \item \textbf{CutPaste}: ritaglia regioni di un'immagine normale, le modifica e le reinserisce per costruire un compito di riconoscimento delle alterazioni;
    \item \textbf{DRAEM}: combina una rete ricostruttiva con una rete discriminativa per apprendere la localizzazione di anomalie simulate \cite{Zavrtanik_et_al._2021};
    \item \textbf{NSA} (\textit{Natural Synthetic Anomalies}): fonde patch provenienti da altre immagini attraverso tecniche di \textit{Poisson image editing} \cite{Cui_et_al._2023}.
\end{itemize}

La difficoltà principale riguarda la relazione tra difetti artificiali e reali: un modello può apprendere caratteristiche della procedura di sintesi che non siano rappresentative delle anomalie industriali. Questi metodi completano il quadro della letteratura, ma non sono inclusi nel confronto sperimentale.

\subsection{Sintesi e collegamento con i modelli selezionati}

La Tabella~\ref{tab:metodi_ad} riassume i principi e i limiti delle famiglie discusse. Le indicazioni sono qualitative: non costituiscono una graduatoria di accuratezza, memoria o velocità indipendente dall'implementazione.

\begin{table}[htbp]
    \centering
    \small
    \caption{Principi e principali limiti delle famiglie di metodi.}
    \label{tab:metodi_ad}
    \renewcommand{\arraystretch}{1.15}
    \begin{tabularx}{\textwidth}{@{}P{0.21\textwidth}YY@{}}
        \hline
        \textbf{Famiglia} & \textbf{Principio ed esempi} & \textbf{Limiti da considerare} \\
        \hline
        Descrittori manuali & Feature progettate e modelli statistici; SIFT, GMM, OCSVM & Scelta delle feature e sensibilità alle condizioni di acquisizione \\
        Ricostruzione & Discrepanza tra input e ricostruzione; autoencoder, CCW-SSIM & Possibile ricostruzione anche dei difetti \\
        Rappresentazioni & Confronto con feature normali o statistiche; PaDiM, PatchCore & Dipendenza dall'estrattore e costo di memoria \\
        Distillazione & Discrepanza Teacher--Student; RD4AD & Dipendenza dalle feature e dalla capacità dello Student \\
        Flussi di normalizzazione & Trasformazione invertibile delle feature; FastFlow & Scelte dell'estrattore, dei flussi e del punteggio \\
        Sintesi di anomalie & Apprendimento da alterazioni artificiali; CutPaste, DRAEM & Corrispondenza tra anomalie sintetiche e reali \\
        \hline
    \end{tabularx}
\end{table}

I quattro modelli selezionati condividono l'adattamento sui soli campioni normali, ma differiscono nello spazio in cui descrivono la normalità e nel criterio con cui ne misurano le deviazioni. Queste differenze motivano un confronto che consideri insieme discriminazione globale, localizzazione e comportamento delle mappe.

\section{Sfide aperte e limiti di applicazione}

Le prestazioni su un benchmark non esauriscono i requisiti di un sistema industriale. Il trasferimento richiede di considerare acquisizione dei dati, cambiamenti del dominio e risorse disponibili \cite{Alzarooni_et_al._2024,Li_et_al._2025}. Le questioni seguenti costituiscono un quadro interpretativo, da distinguere dalle proprietà effettivamente misurate nel progetto.

\subsection{Differenze di dominio e adattamento delle feature}

I backbone pre-addestrati su dataset generici, come ImageNet, rendono disponibili rappresentazioni senza richiedere di apprendere da zero l'intero estrattore \cite{Roth_et_al._2022}. Tuttavia, le caratteristiche utili alla classificazione di immagini naturali possono non descrivere in modo ottimale microstrutture e difetti industriali.

L'adattamento delle feature può ridurre questa distanza, ma richiede di verificare che la rappresentazione conservi la capacità di distinguere anomalie non osservate \cite{Li_et_al._2025}. Nel progetto, PaDiM, RD4AD e FastFlow condividono un backbone congelato; CCW-SSIM apprende invece dai soli campioni industriali. Questa differenza deve essere considerata nell'interpretazione dei risultati.

\subsection{Scalabilità computazionale ed edge computing}

L'esecuzione su dispositivi periferici può imporre vincoli di memoria, latenza ed energia, discussi sia per l'ispezione visiva industriale sia, più in generale, per le reti IoT \cite{Li_et_al._2025,Morshedi_et_al._2025}. Il costo dipende dall'estrattore, dalla risoluzione, dai componenti addestrabili e dal trattamento delle mappe.

Nei metodi con memory bank, la selezione di un coreset limita la quantità di feature conservate, ma la consultazione dei riferimenti rimane parte dell'inferenza \cite{Roth_et_al._2022}. Nei modelli neurali, il costo dell'addestramento deve essere distinto dal tempo di predizione. Per questo, un tempo di addestramento basso non dimostra da solo l'idoneità a operare in tempo reale.

\subsection{Cambiamenti del dominio e apprendimento continuo}

Molti esperimenti assumono che le condizioni di addestramento e test siano stabili. Nelle applicazioni possono invece cambiare illuminazione, strumenti di acquisizione o caratteristiche del prodotto. Il modello appreso può così perdere validità e produrre più falsi allarmi.

La \textit{continual anomaly detection} studia aggiornamenti incrementali della normalità, affrontando il rischio di dimenticare informazioni precedenti (\textit{catastrophic forgetting}) o di incorporare anomalie non riconosciute. Una rassegna recente di questi temi è disponibile come preprint, non sottoposto a peer review \cite{Faber_et_al._2026}. La presenza di anomalie mai viste va distinta da questi cambiamenti della normalità: sono due aspetti diversi della generalizzazione.

\subsection{Approcci zero-shot e foundation models}

Quando è difficile raccogliere campioni normali specifici del prodotto, gli approcci \textit{zero-shot} cercano di rilevare anomalie senza adattamento sul dominio di destinazione. Modelli visivo-linguistici e altri foundation models possono fornire rappresentazioni e indicazioni semantiche mediante prompt.

Il framework DCS esplora, ad esempio, l'integrazione di DINO, CLIP e SAM per il rilevamento e la segmentazione \cite{Wan_et_al._2026}. Queste possibilità ampliano il quadro della ricerca, ma richiedono valutazioni specifiche della precisione locale e dei costi. Non fanno parte dei metodi sperimentati nel presente lavoro.

La rassegna conduce quindi a una domanda circoscritta: come si comportano quattro diverse pipeline one-class quando sono adattate e valutate sullo stesso caso di studio? Il capitolo seguente documenta le configurazioni necessarie a rendere esplicite le condizioni di tale confronto.
